% =============================================================================
% BLOC 1 — 20 minutes, dont 3 à 5 minutes en anglais à la fin
% =============================================================================
\section{Bloc 1 — Planifier et organiser}
\sectiondivider{Bloc 1}{Planifier et organiser un projet de développement logiciel}{Fil rouge : Lynx Immo · Ancrage professionnel : FCU / Ediacara}

\begin{frame}{Bloc 1 — De l'opportunité au projet piloté}
\context{B1}{C1.1 $\rightarrow$ C1.7}
\begin{columns}[T]
\begin{column}{0.50\linewidth}
\begin{block}{Lynx Immo : pilotage structuré}
\begin{tightitems}
  \item C1.1 — Faisabilité.
  \item C1.2 — Veille technologique.
  \item C1.4 — Feuille de route.
  \item C1.5 — Mise en place.
  \item C1.6 — Mesure de la performance.
  \item C1.7 — Supervision et arbitrages.
\end{tightitems}
\end{block}
\end{column}
\begin{column}{0.46\linewidth}
\begin{block}{FCU / Ediacara : décision en contexte réel}
\begin{tightitems}
  \item C1.3 — Concevoir une adaptation compatible avec une production historique.
  \item Arbitrer entre refonte idéale et évolution maîtrisée.
  \item Tenir compte du risque opérationnel et du legacy.
\end{tightitems}
\end{block}
\end{column}
\end{columns}
\end{frame}

\begin{frame}{Lynx Immo — Définir un produit suffisamment concret pour piloter}
\context{B1}{Support projet}
\begin{columns}[T]
\begin{column}{0.46\linewidth}
\begin{tightitems}
  \item Client réel : organisme de formation en immobilier commercial.
  \item Besoin initial : micro-apprentissage mobile.
  \item Produit élargi : application apprenant, administration, back-end, abonnements et données d'usage.
  \item Mon rôle : \strong{chef de projet et référent technique} au sein d'une équipe de cinq étudiants.
\end{tightitems}
\end{column}
\begin{column}{0.50\linewidth}
\centering
\repofigure[0.48\textheight]{figures/ipf/priorisation_perimetre.pdf}
\end{column}
\end{columns}
\end{frame}

\begin{frame}{C1.1 — Vérifier la faisabilité avant de s'engager}
\context{B1}{C1.1 · Lynx Immo}
\begin{block}{Question de départ}
Le produit souhaité est-il réalisable dans le temps, avec l'équipe et les ressources disponibles ?
\end{block}
\begin{columns}[T]
\begin{column}{0.49\linewidth}
\begin{tightitems}
  \item Définition d'un \strong{MVP} avant la distribution des tâches.
  \item Comparaison de scénarios d'architecture et de leurs coûts d'exploitation.
  \item Analyse des dépendances externes : authentification, paiements, stores et hébergement.
  \item Identification des risques techniques, humains et financiers.
\end{tightitems}
\end{column}
\begin{column}{0.47\linewidth}
\begin{block}{Décision}
Ne pas traiter la faisabilité comme une validation binaire, mais comme un \textbf{outil de cadrage du périmètre}.
\end{block}
\vspace{0.5em}
\small
\strong{Résultat :} un produit priorisé, des fonctions différées et une architecture compatible avec le contexte réel du projet.
\end{column}
\end{columns}
\end{frame}

\begin{frame}{C1.3 — Concevoir dans les contraintes d'une chaîne legacy}
\context{B1}{C1.3 · FCU / Ediacara}
\begin{columns}[T]
\begin{column}{0.45\linewidth}
\begin{tightitems}
  \item Le flux historique part de la BDMR ; le FCU n'y passe jamais.
  \item Les deux flux doivent pourtant atteindre les \strong{mêmes espaces de diffusion}.
  \item Risque principal : une synchronisation historique peut considérer un produit FCU comme orphelin et le supprimer.
  \item Les exécutables Eiffel sont partagés avec d'autres briques et ne peuvent pas être refondus sans risque majeur.
\end{tightitems}
\end{column}
\begin{column}{0.51\linewidth}
\centering
\repofigure[0.50\textheight]{figures/fcu/coexistence_flux_historique_fcu.png}
\end{column}
\end{columns}
\begin{block}{Décision d'architecture}
\textbf{Séparer explicitement le flux FCU tout en sécurisant tous les points de contact avec le flux historique.}
\end{block}
\end{frame}

\begin{frame}{C1.4 — Transformer un produit large en trajectoire}
\context{B1}{C1.4 · Lynx Immo}
\begin{columns}[T]
\begin{column}{0.40\linewidth}
\begin{tightitems}
  \item Feuille de route de novembre à juillet.
  \item Jalons associés à des livrables et des décisions.
  \item Matrice d'exigences pour garder le lien avec le besoin client.
  \item Méthode hybride plutôt qu'un Scrum appliqué mécaniquement.
\end{tightitems}
\begin{block}{Usage réel}
La roadmap a servi à \textbf{arbitrer le périmètre, les priorités et la charge}, pas seulement à afficher des dates.
\end{block}
\end{column}
\begin{column}{0.56\linewidth}
\centering
\repofigure[0.52\textheight]{figures/ipf/roadmap_calendrier.pdf}
\end{column}
\end{columns}
\end{frame}

\begin{frame}{C1.5 — Mettre en place un environnement de travail commun}
\context{B1}{C1.5 · Lynx Immo}
\begin{columns}[T]
\begin{column}{0.48\linewidth}
\begin{block}{Ce que j'ai mis en place}
\begin{tightitems}
  \item Découpage du produit en tâches manipulables.
  \item Jira pour le suivi opérationnel.
  \item Documentation des décisions et exigences.
  \item Répartition par domaines techniques et fonctionnels.
\end{tightitems}
\end{block}
\end{column}
\begin{column}{0.48\linewidth}
\begin{block}{Ce que j'ai appris}
\begin{tightitems}
  \item La capacité théorique n'est pas la capacité réelle.
  \item Le temps d'apprentissage fait partie de la charge.
  \item Une méthode doit s'adapter à l'équipe plutôt que l'inverse.
\end{tightitems}
\end{block}
\end{column}
\end{columns}
\vspace{0.5em}
\centering
\repofigure[0.30\textheight]{figures/ipf/methode_hybride_5sc.png}
\end{frame}

\begin{frame}{C1.6 — Mesurer l'avancement sans confondre activité et résultat}
\context{B1}{C1.6 · Lynx Immo}
\begin{columns}[T]
\begin{column}{0.48\linewidth}
\begin{block}{Quatre familles de KPI}
\begin{tightitems}
  \item \strong{Projet} : jalons, tâches, exigences.
  \item \strong{Valeur / usage} : engagement, continuité d'apprentissage.
  \item \strong{Technique} : temps de réponse, disponibilité, erreurs.
  \item \strong{Charge} : comportement sous sollicitation contrôlée.
\end{tightitems}
\end{block}
\end{column}
\begin{column}{0.48\linewidth}
\centering
\repofigure[0.43\textheight]{figures/ipf/vue_synthetique_avancement_jira.png}
\end{column}
\end{columns}
\begin{block}{Principe}
Un indicateur n'a de valeur que s'il \textbf{déclenche ou éclaire une décision}.
\end{block}
\end{frame}

\begin{frame}{C1.7 — Superviser, arbitrer et réajuster}
\context{B1}{C1.7 · Lynx Immo}
\begin{columns}[T]
\begin{column}{0.46\linewidth}
\begin{tightitems}
  \item Suivi des écarts entre planning et capacité réelle.
  \item Registre des risques avec criticité et règle d'escalade.
  \item Arbitrages avec le client sans perdre la trajectoire du projet.
  \item Réallocation de tâches et réduction de périmètre lorsqu'un risque devient concret.
\end{tightitems}
\begin{block}{Logique de supervision}
\textbf{Constater $\rightarrow$ analyser $\rightarrow$ décider $\rightarrow$ mesurer l'effet de l'ajustement.}
\end{block}
\end{column}
\begin{column}{0.50\linewidth}
\centering
\repofigure[0.50\textheight]{figures/ipf/matrice_risque-5sc.pdf}
\end{column}
\end{columns}
\end{frame}

% ----------------------------------------------------------------------------
% ENGLISH PART — target 4 min
% ----------------------------------------------------------------------------
\begin{frame}{C1.2 — Technology watch as a decision-making system}
\selectlanguage{english}
\context{B1 · ENGLISH}{C1.2 · Lynx Immo}
\begin{block}{Objective}
The goal was not to collect technology news. It was to build a \textbf{repeatable decision process} for architecture choices.
\end{block}
\begin{tightitems}
  \item Define the scope of the watch and the decisions it must support.
  \item Prioritize official English-language documentation and primary sources.
  \item Compare options using explicit criteria: maturity, maintainability, cost, team skills and deployment constraints.
  \item Define triggers that require the decision to be reviewed when assumptions change.
\end{tightitems}
\begin{block}{Transversal competency}
The written watch report is in French, but it is based on English technical sources and required me to extract, compare and explain technical information.
\end{block}
\end{frame}

\begin{frame}{Replaying the decision when assumptions changed}
\context{B1 · ENGLISH}{C1.2 · Lynx Immo}
\begin{columns}[T]
\begin{column}{0.45\linewidth}
\begin{block}{Initial assumptions}
\begin{tightitems}
  \item React Native for the mobile application.
  \item Django for the back-end.
  \item More infrastructure to operate ourselves.
\end{tightitems}
\end{block}
\begin{block}{New information}
The real product scope, team skills, delivery constraints and service needs changed the evaluation.
\end{block}
\end{column}
\begin{column}{0.51\linewidth}
\begin{block}{Final architecture}
\begin{tightitems}
  \item Flutter for the learner application.
  \item React for the administration interface.
  \item Supabase / PostgreSQL for the back-end platform.
  \item RevenueCat for subscription management.
\end{tightitems}
\end{block}
\centering
\repofigure[0.22\textheight]{figures/ipf/architecture.png}
\end{column}
\end{columns}
\vspace{0.25em}\small\strong{Key lesson:} A good technology watch must help us \textbf{change our mind for documented reasons}.
\selectlanguage{french}
\end{frame}

\begin{frame}{Bloc 1 — Synthèse des preuves}
\context{B1}{Synthèse}
\small
\begin{tabularx}{\linewidth}{>{\bfseries}p{1.35cm}p{4.2cm}X}
\toprule
Comp. & Projet principal & Preuve retenue \\
\midrule
C1.1 & Lynx Immo & Étude de faisabilité et cadrage du MVP. \\
C1.2 & Lynx Immo & Système de veille et réévaluation des choix. \\
C1.3 & FCU / Ediacara & Solution d'intégration compatible avec le legacy. \\
C1.4 & Lynx Immo & Roadmap, jalons et exigences. \\
C1.5 & Lynx Immo & Environnement de travail et méthode. \\
C1.6 & Lynx Immo & KPI projet, usage, technique et charge. \\
C1.7 & Lynx Immo & Risques, arbitrages et réajustements. \\
\bottomrule
\end{tabularx}
\end{frame}
